
\section{\BookRecordIndex{T0327}椭圆估计}\label{sec:C-2}

本节将证明\BookRecordIndex{T0210}柯西–黎曼算子的弗雷德霍姆性质以及定理~\ref{thm:C-1-10}(i) 中的对偶断言。主要关注指标公式证明的读者，可以直接读 C.3 和 C.4 节。

弗雷德霍姆性质的证明以引理~\ref{lem:A-1-1} 为基础：若算子 $D$ 满足某个估计，则它的像闭合，核有限维。光滑\BookRecordIndex{T0210}柯西–黎曼算子所需的估计由下述引理~\ref{lem:C-2-1} 建立，因而在光滑情形，$D_F$ 和 $D_F^*$ 都有有限维核与闭像。一般情形也一样，因为如下引理~\ref{lem:C-2-2} 将证明，每个 $W^{\ell-1,p}$ 类\BookRecordIndex{T0210}柯西–黎曼算子都是某个光滑\BookRecordIndex{T0210}柯西–黎曼算子的紧扰动。还须证明余核也有限维；这来自一个对偶定理，它把 $D_F$ 的余核与 $D_F^*$ 的核相互等同。下述定理~\ref{thm:C-2-3} 给出这一对偶断言的加强形式。证明以\BookRecordIndex{T0329}椭圆正则性为基础，由于 $D$ 的系数只具较弱的\BookRecordIndex{T0408}正则性，需要分成若干步骤。不过，如步骤 2、3 所说明，在光滑情形，结论直接来自引理~\ref{lem:C-1-9} 与 B.4 节的\BookRecordIndex{T0408}正则性结果。

以下始终假定 $(\Sigma,j,\mathrm{dvol})$ 是带边紧致黎曼曲面，$E$ 是 $\Sigma$ 上的复向量丛，$F\subset E|_{\partial\Sigma}$ 是全实子丛，$\langle\cdot,\cdot\rangle$ 是 $E$ 上满足 $JF=F^\perp$ 的埃尔米特内积。

\begin{lemma}\label{lem:C-2-1}
设 $D$ 是光滑的实线性\BookRecordIndex{T0210}柯西–黎曼算子。则对每个正整数 $k$ 及每个 $q>1$，存在常数 $c>0$，使对每个 $\xi\in W_F^{k,q}(\Sigma,E)$，
\[
  \|\xi\|_{W^{k,q}}\le c\bigl(\|D_F\xi\|_{W^{k-1,q}}+\|\xi\|_{W^{k-1,q}}\bigr),
\]
且对每个 $\eta\in W_F^{k,q}(\Sigma,\Lambda^{0,1}T^*\Sigma\otimes_{\mathbb C}E)$，
\[
  \|\eta\|_{W^{k,q}}\le c\bigl(\|D_F^*\eta\|_{W^{k-1,q}}+\|\eta\|_{W^{k-1,q}}\bigr).
\]
\end{lemma}
\begin{proof}
在 $\Sigma$ 上的局部全纯坐标图中取 $E$ 的局部平凡化，方程 $D_F\xi=\zeta$ 写作
\begin{equation}\label{eq:C-2-1}
  \partial_s\xi+\Phi\xi+i(\partial_t\xi+\Psi\xi)=\zeta,
  \qquad \xi(s,0)\in\mathbb R^n.
\end{equation}
这里 $s+it$ 是上半平面
\[
  \mathbb H:=\{s+it\in\mathbb C\mid t\ge0\}
\]
的开子集 $\Omega$ 上的坐标，$\Phi\,ds+\Psi\,dt$ 是相应联络 $\nabla$ 的\BookRecordIndex{T0223}联络势（见习题~\ref{ex:C-1-6}），故 $\Phi,\Psi\in C^\infty(\Omega,\operatorname{End}_{\mathbb R}(\mathbb C^n))$。函数 $\zeta\in W^{k-1,q}(\Omega,\mathbb C^n)$ 表示 $(0,1)$-形式 $\frac12(\zeta\,ds-i\zeta\,dt)$。对任意闭包包含于 $\Omega$ 的开集 $\Omega'$，命题~\ref{prop:B-4-9} 给出正数 $c,c'$，使方程~\eqref{eq:C-2-1} 的每个解 $\xi\in W^{k,q}_{\mathrm{loc}}(\Omega,\mathbb C^n)$ 满足
\[
\begin{aligned}
  \|\xi\|_{W^{k,q}(\Omega')}
  &\le c\bigl(\|\zeta-\Phi\xi-J\Psi\xi\|_{W^{k-1,q}(\Omega)}
                  +\|\xi\|_{W^{k-1,q}(\Omega)}\bigr)\\
  &\le c'\bigl(\|\zeta\|_{W^{k-1,q}(\Omega)}
                  +\|\xi\|_{W^{k-1,q}(\Omega)}\bigr).
\end{aligned}
\]
为得到全局估计，用有限个与 $\mathbb H$ 中开球双全纯的开集覆盖 $\Sigma$，并在每个开集上取 $E$ 的\BookRecordIndex{T0391}酉平凡化。第一项得证。

现在用 $E^*$ 表示同一底\BookRecordIndex{T0072}层丛 $E$，但将复结构换成 $-J$。由注~\ref{rem:C-1-3}，$D^*$ 与丛 $T\Sigma\otimes_{\mathbb C}E^*$ 上、边界全实子丛为 $T\partial\Sigma\otimes_{\mathbb R}F$ 的\BookRecordIndex{T0210}柯西–黎曼算子共轭。因此第二项由第一项\BookRecordIndex{T0326}推出。
\end{proof}

由引理~\ref{lem:A-1-1} 和~\ref{lem:C-2-1} 可知，在光滑情形，$D_F$、$D_F^*$ 的核均有限维且像闭合。推论~\ref{cor:A-1-2} 与下一引理说明，此结论也适用于 $W^{\ell-1,p}$ 类\BookRecordIndex{T0210}柯西–黎曼算子。

\begin{lemma}\label{lem:C-2-2}
设 $\ell$ 为整数，$p>1$ 且 $\ell p>2$，并设
\[
  A\in W^{\ell-1,p}(\Sigma,T^*\Sigma\otimes_{\mathbb R}\operatorname{End}_{\mathbb R}(E)).
\]
则对任意 $1\le k\le\ell$、$q>1$ 且 $k-2/q\le\ell-2/p$，线性算子
\[
  W^{k,q}(\Sigma,E)\longrightarrow
  W^{k-1,q}(\Sigma,\Lambda^{0,1}T^*\Sigma\otimes E):
  \xi\longmapsto(A\xi)^{0,1}
\]
良定义且紧。
\end{lemma}
\begin{proof}
若 $k=\ell$，则 $q\le p$，由 Hölder 不等式有 $W^{\ell-1,p}\subset W^{k-1,q}$。若 $k\le\ell-2$，则定理~\ref{thm:B-1-11} 给出对所有 $q$ 的包含
\[
  W^{\ell-1,p}\subset C^{k-1}\subset W^{k-1,q}.
\]
若 $k=\ell-1$，不等式 $k-2/q\le\ell-2/p$ 等价于 $2/p-1\le2/q$；因此在 $p<2$ 时有 $q\le2p/(2-p)$。于是由定理~\ref{thm:B-1-12}，仍有 $W^{\ell-1,p}\subset W^{k-1,q}$。综上，$A\in W^{k-1,q}$。

先设 $kq>2$。映射 $W^{k,q}\to W^{k-1,q}$，$\xi\mapsto(A\xi)^{0,1}$，分解经过包含映射 $W^{k,q}\hookrightarrow C^{k-1}$。由定理~\ref{thm:B-1-11}，该包含映射是紧的，故引理在 $kq>2$ 时成立。若 $kq\le2$，则 $k=1$、$q\le2$。此外，由定理~\ref{thm:B-1-12}，存在 $s>2$ 使 $A\in L^s$。取
\[
  r:=\frac{sq}{s-q}<\frac{2q}{2-q}.
\]
由定理~\ref{thm:B-1-12}，包含 $W^{1,q}\hookrightarrow L^r$ 是紧的\sourcefootnote{原书上式在 $q=2$ 时的上界 $2q/(2-q)$ 分母为零。此时 $r=2s/(s-2)<\infty$；取充分接近 $2$ 的 $q'<2$，由 $W^{1,2}\hookrightarrow W^{1,q'}$ 及定理 B.1.12，仍得到 $W^{1,2}\hookrightarrow L^r$ 的\BookRecordIndex{T0192}紧性。}。又因 $1/r+1/s=1/q$，Hölder 不等式表明乘以 $A$ 给出从 $L^r$ 到 $L^q$ 的有界线性算子。故引理在 $kq\le2$ 时也成立。
\end{proof}

为了证明\BookRecordIndex{T0210}柯西–黎曼算子的弗雷德霍姆性质，须将 $D_F$ 的余核与 $D_F^*$ 的核等同。这正是下述重述定理~\ref{thm:C-1-10}(i) 的定理所论述的。

\begin{theorem}\label{thm:C-2-3}
设 $D$ 是 $W^{\ell-1,p}$ 类的实线性\BookRecordIndex{T0210}柯西–黎曼算子，其中 $\ell$ 为正整数，$p>1$ 且 $\ell p>2$。设整数 $k$ 满足 $1\le k\le\ell$，$q,r>1$ 满足 $2/r-1\le k-2/q\le\ell-2/p$。则：
\begin{enumerate}[label=(\roman*)]
\item 算子
\[
  D_F:W_F^{k,q}(\Sigma,E)\to W^{k-1,q}(\Sigma,\Lambda^{0,1}T^*\Sigma\otimes_{\mathbb C}E),
\]
\[
  D_F^*:W_F^{k,q}(\Sigma,\Lambda^{0,1}T^*\Sigma\otimes_{\mathbb C}E)
  \to W^{k-1,q}(\Sigma,E)
\]
为\BookRecordIndex{T0128}弗雷德霍姆算子，且 $\operatorname{index}(D_F)+\operatorname{index}(D_F^*)=0$。$D_F$ 的\BookRecordIndex{T0129}弗雷德霍姆指标与 $D$ 的选择、$\Sigma$ 上的复结构以及 $k,q$ 的选择都无关。
\item 若 $\eta\in L^r(\Sigma,\Lambda^{0,1}T^*\Sigma\otimes_{\mathbb C}E)$ 和 $\xi\in W^{k-1,q}(\Sigma,E)$ 满足
\begin{equation}\label{eq:C-2-2}
  \int_\Sigma\langle\eta,D_F\zeta\rangle\,\mathrm{dvol}
  =\int_\Sigma\langle\xi,\zeta\rangle\,\mathrm{dvol}
\end{equation}
对所有 $\zeta\in W_F^{k,q}(\Sigma,E)$ 成立，则 $\eta\in W_F^{k,q}(\Sigma,\Lambda^{0,1}T^*\Sigma\otimes_{\mathbb C}E)$，且 $D_F^*\eta=\xi$。
\item 若 $\xi\in L^r(\Sigma,E)$、$\eta\in W^{k-1,q}(\Sigma,\Lambda^{0,1}T^*\Sigma\otimes_{\mathbb C}E)$ 满足
\begin{equation}\label{eq:C-2-3}
  \int_\Sigma\langle D_F^*\zeta,\xi\rangle\,\mathrm{dvol}
  =\int_\Sigma\langle\zeta,\eta\rangle\,\mathrm{dvol}
\end{equation}
对所有 $\zeta\in W_F^{k,q}(\Sigma,\Lambda^{0,1}T^*\Sigma\otimes_{\mathbb C}E)$ 成立，则 $\xi\in W_F^{k,q}(\Sigma,E)$，且 $D_F\xi=\eta$。
\end{enumerate}
\end{theorem}
\begin{proof}
证明分八步。步骤 1 说明 (ii)、(iii) 中的假设有意义；步骤 2、3 对光滑\BookRecordIndex{T0210}柯西–黎曼算子证明结论。

\paragraph{步骤 1。}（a）取 $s>1$ 使 $1/s+1/r=1$。则有包含 $W^{k-1,q}\hookrightarrow L^s$ 及 $L^r\hookrightarrow(W^{k-1,q})^*$。

（b）有包含 $W^{k,q}\hookrightarrow(W^{\ell-1,p})^*$。

（c）当且仅当 $-1/2\le k-2/q$ 时，有包含 $W^{k,q}\hookrightarrow(W^{k-1,q})^*$。

不等式 $-2/s=2/r-2\le k-1-2/q$ 表明：若 $(k-1)q<2$，则 $s\le2q/(2-(k-1)q)$。因此，$W^{k-1,q}\hookrightarrow L^s$ 由\BookRecordIndex{T0318}索伯列夫嵌入定理~\ref{thm:B-1-12} 给出，证明了 (a)。

若 $kq\ge2$ 或 $(\ell-1)p\ge2$，(b) 显然。故设 $kq<2$ 且 $(\ell-1)p<2$。定理~\ref{thm:B-1-12} 给出 $W^{k,q}\hookrightarrow L^{2q/(2-kq)}$，且又给出稠密包含
\[
  W^{\ell-1,p}\hookrightarrow L^{2p/(2+p-\ell p)},
  \qquad
  L^{2p/(p+\ell p-2)}\hookrightarrow(W^{\ell-1,p})^*.
\]
因为 $\ell p>2$，有 $2/q-k\le1\le1+\ell-2/p$，从而 $2q/(2-kq)\ge2p/(p+\ell p-2)$。这蕴涵
\[
  W^{k,q}\subset L^{2q/(2-kq)}
  \subset L^{2p/(p+\ell p-2)}
  \subset(W^{\ell-1,p})^*.
\]
故 (b) 得证。

当 $kq\ge2$ 时 (c) 显然。设 $kq<2$，于是 $k=1$、$1<q<2$。由\BookRecordIndex{T0318}索伯列夫嵌入定理~\ref{thm:B-1-12}，
\[
  W^{1,q}\hookrightarrow(L^q)^*\cong L^{q/(q-1)}
  \quad\Longleftrightarrow\quad
  \frac{q}{q-1}\le\frac{2q}{2-q}.
\]
这当且仅当 $q\ge4/3$，等价于 $-1/2\le1-2/q=k-2/q$。步骤 1 得证。

\paragraph{步骤 2。} 假定 $D$ 是光滑\BookRecordIndex{T0210}柯西–黎曼算子，证明 (ii) 和 (iii)。

考虑 $\Sigma$ 上局部全纯坐标 $s+it$ 和局部平凡化下的弱方程，其中体积形式为 $\lambda^2ds\wedge dt$。此时式~\eqref{eq:C-2-2} 写作
\begin{equation}\label{eq:C-2-4}
  \int_\Omega\langle\eta,\partial_s\zeta+i\partial_t\zeta\rangle\,ds\,dt
  =\int_\Omega\langle\lambda^2\xi-\Phi^*\eta+\Psi^*i\eta,\zeta\rangle\,ds\,dt,
\end{equation}
其中 $\eta\in L^r_{\mathrm{loc}}(\Omega,\mathbb C^n)$ 表示 $(0,1)$-形式 $\frac12(\eta\,ds-i\eta\,dt)$，$\xi\in W^{k-1,q}_{\mathrm{loc}}(\Omega,\mathbb C^n)$；测试函数 $\zeta\in C_0^\infty(\Omega,\mathbb C^n)$ 对 $\Omega\cap\mathbb R$ 上的 $s$ 满足 $\zeta(s,0)\in\mathbb R^n$。依假设，$\Phi,\Psi:\Omega\to\operatorname{End}_{\mathbb R}(\mathbb C^n)$ 光滑，$\Phi^*,\Psi^*$ 表示其（实）伴随矩阵。设 $\eta$ 对所有满足 $\zeta(\Omega\cap\mathbb R)\subset\mathbb R^n$ 的上述测试函数都满足式~\eqref{eq:C-2-4}。须证明 $\eta\in W^{k,q}_{\mathrm{loc}}(\Omega,\mathbb C^n)$。为此，注意
\[
  \lambda^2\xi-\Phi^*\eta+\Psi^*i\eta
  \in L^{r_0}_{\mathrm{loc}}(\Omega,\mathbb C^n),
  \qquad r_0:=\min\{r,q\}>1.
\]
故由引理~\ref{lem:B-4-6}（取 $p=\infty$），$\eta\in W^{1,r_0}_{\mathrm{loc}}$，且
\[
  -\partial_s\eta+i\partial_t\eta
  =\lambda^2\xi-\Phi^*\eta+\Psi^*i\eta,
  \qquad \eta(s,0)\in\mathbb R^n.
\]
若 $r_0\ge2$，定理~\ref{thm:B-1-12} 给出 $\eta\in L^q_{\mathrm{loc}}$。若 $r_0<2$，则 $\eta\in L^{2r_0/(2-r_0)}_{\mathrm{loc}}$，因而
\[
  \lambda^2\xi-\Phi^*\eta+\Psi^*i\eta
  \in L^{r_1}_{\mathrm{loc}}(\Omega,\mathbb C^n),
  \qquad r_1:=\min\{2r_0/(2-r_0),q\}.
\]
数 $r_1$ 或等于 $q$，或大于 $2$。后一情形再由引理~\ref{lem:B-4-6} 得 $\eta\in W^{1,2}_{\mathrm{loc}}$，进而由定理~\ref{thm:B-1-12} 得 $\eta\in L^q_{\mathrm{loc}}$。现由命题~\ref{prop:B-4-9} 可知：若对某个 $j\in\{1,\ldots,k\}$ 已有 $\eta\in W^{j-1,q}_{\mathrm{loc}}$，则 $\eta\in W^{j,q}_{\mathrm{loc}}$。归纳得到 $\eta\in W^{k,q}_{\mathrm{loc}}$。这里命题~\ref{prop:B-4-9} 中的 $p>2$ 条件是多余的，因为目前的（殆）复结构光滑（见习题~\ref{ex:B-4-11}）。

为完成步骤 2，用有限多个全纯坐标图覆盖 $\Sigma$，并在这些图上取 $E$ 的\BookRecordIndex{T0391}酉平凡化，应用刚得出的局部结果。

\paragraph{步骤 3。} 假定 $D$ 光滑，证明 (i)。

先设 $k=1$、$1/r+1/q=1$。则 $2/r-1=1-2/q$，所以 $r$ 满足定理中 $k=1$ 的假设。将 $D_F$ 与 $D_F^*$ 视为 $W_F^{1,q}\to L^q$ 的算子，我们断言
\begin{equation}\label{eq:C-2-5}
  L^q(\Sigma,\Lambda^{0,1}T^*\Sigma\otimes_{\mathbb C}E)
  =\operatorname{im}D_F\oplus\ker D_F^*.
\end{equation}
先证明交为零。若 $\eta_0\in\operatorname{im}D_F\cap\ker D_F^*$，则 $\eta_0\in W_F^{1,q}(\Sigma,\Lambda^{0,1}T^*\Sigma\otimes_{\mathbb C}E)$。由于 $D_F^*\eta_0=0$，引理~\ref{lem:C-1-9} 表明 $\eta_0$ 湮灭 $D_F:W_F^{1,r}\to L^r$ 的像。由步骤 2 交换 $r,q$ 的角色，$\eta_0$ 光滑。\footnote{此处可清楚看到证明的主要难点。若 $k=\ell=1$、$q=p$ 且 $1/r+1/q=1$，当 $D$ 仅为 $L^p$ 类\BookRecordIndex{T0210}柯西–黎曼算子时，我们只知道 $\lambda^2\xi-\Phi^*\eta+\Psi^*i\eta$ 局部可积，却没有 $L^1$ \BookRecordIndex{T0329}椭圆正则性理论；\BookRecordIndex{T0012}Calderón--Zygmund 不等式只对 $p>1$ 成立。主要应用中正会出现这种情况。}（注意，$k-2/r\le\ell-2/p$ 未必成立，但对光滑算子而言，这一条件无关紧要。）因此再由引理~\ref{lem:C-1-9}，$\eta_0$ 也湮灭 $D_F:W_F^{1,q}\to L^q$ 的像。由于 $\eta_0$ 属于该像，$\|\eta_0\|_{L^2}^2=0$，故 $\eta_0=0$。

还要证明直和为整个 $L^q(\Sigma,\Lambda^{0,1}T^*\Sigma\otimes_{\mathbb C}E)$。由引理~\ref{lem:C-2-1}，$D_F$ 的像闭合、$D_F^*$ 的核有限维，故其直和是 $L^q$ 的闭子空间。只需证明它稠密；由 Hahn--Banach 定理，只需其湮灭子为零。设 $\zeta\in L^r(\Sigma,\Lambda^{0,1}T^*\Sigma\otimes_{\mathbb C}E)$ 满足
\[
  \zeta\perp\operatorname{im}D_F\oplus\ker D_F^*.
\]
由 $\zeta\perp\operatorname{im}D_F$ 和步骤 2，$\zeta$ 光滑且 $D_F^*\zeta=0$。再由 $\zeta\perp\ker D_F^*$ 得 $\zeta=0$。这证明式~\eqref{eq:C-2-5}。交换 $D_F$、$D_F^*$ 的角色也得到类似公式，故二者都是\BookRecordIndex{T0128}弗雷德霍姆算子，且 $\operatorname{index}(D_F)+\operatorname{index}(D_F^*)=0$。这证明 $k=1$ 的断言。

现在设 $k\ge2$，并将 $D_F,D_F^*$ 视为 $W_F^{k,q}\to W^{k-1,q}$ 的算子。我们断言
\begin{equation}\label{eq:C-2-6}
\begin{aligned}
  W^{k-1,q}(\Sigma,\Lambda^{0,1}T^*\Sigma\otimes_{\mathbb C}E)
  &=\operatorname{im}D_F\oplus\ker D_F^*,\\
  W^{k-1,q}(\Sigma,E)
  &=\operatorname{im}D_F^*\oplus\ker D_F.
\end{aligned}
\end{equation}
先证第一式。交为零，因为 $D_F^*$ 的核由光滑截面组成，而 $D_F$ 的像包含在 $D_F:W_F^{1,q}\to L^q$ 的像中。任取 $\eta\in W^{k-1,q}(\Sigma,\Lambda^{0,1}T^*\Sigma\otimes_{\mathbb C}E)$，由式~\eqref{eq:C-2-5} 可写为
\[
  \eta=D_F\xi+\eta_0,\qquad D_F^*\eta_0=0,
\]
其中 $\eta_0\in W^{1,q}(\Sigma,\Lambda^{0,1}T^*\Sigma\otimes_{\mathbb C}E)$，$\xi\in W_F^{1,q}(\Sigma,E)$。如前，$\eta_0$ 光滑，因此 $D_F\xi=\eta-\eta_0\in W^{k-1,q}$。由\BookRecordIndex{T0318}索伯列夫嵌入定理~\ref{thm:B-1-11} 和~\ref{thm:B-1-12}，$\xi$ 连续。再由已在步骤 2 对光滑\BookRecordIndex{T0210}柯西–黎曼算子证明的 (iii)，得 $\xi\in W_F^{k,q}(\Sigma,E)$，从而式~\eqref{eq:C-2-6} 的第一式成立。交换 $D_F,D_F^*$ 的角色得到第二式。由式~\eqref{eq:C-2-6} 可知二者为\BookRecordIndex{T0128}弗雷德霍姆算子。其指标与 $k,q$ 无关，因为二者的核都是光滑截面。不同\BookRecordIndex{T0210}柯西–黎曼算子的差是从 $W_F^{k,q}(\Sigma,E)$ 到 $W^{k-1,q}(\Sigma,\Lambda^{0,1}T^*\Sigma\otimes E)$ 的紧算子，这是 Rellich 定理的结论（见定理~\ref{thm:B-1-12} 与定理~\ref{thm:A-1-5}）；故 $D_F$ 的指标与 $D$ 的选择无关。

为说明它与 $\Sigma$ 上复结构的选择也无关，沿一族光滑复结构 $\lambda\mapsto j_\lambda$ 作\BookRecordIndex{T0263}平行移动，构造丛同构的光滑族 $\Phi_\lambda:T\Sigma\to T\Sigma$，满足 $\Phi_0=\operatorname{id}$、$\Phi_\lambda j_0=j_\lambda\Phi_\lambda$。在 $E$ 上固定\BookRecordIndex{T0052}埃尔米特联络 $\nabla$；对每个 $\lambda$，令 $D_\lambda$ 为由 $\nabla$ 与 $j_\lambda$ 关联的\BookRecordIndex{T0210}柯西–黎曼算子。于是有连续的\BookRecordIndex{T0128}弗雷德霍姆算子族
\[
  W_F^{k,q}(\Sigma,E)
  \longrightarrow
  W^{k-1,q}(\Sigma,\Lambda^{0,1}_{j_0}T^*\Sigma\otimes E):
  \xi\longmapsto(D_\lambda\xi)\circ\Phi_\lambda.
\]
由定理~\ref{thm:A-1-5}，此族算子及 $D_\lambda$ 的\BookRecordIndex{T0129}弗雷德霍姆指标与 $\lambda$ 无关。步骤 3 得证。

至此已经对光滑\BookRecordIndex{T0210}柯西–黎曼算子证明定理。余下部分将结果延伸到 $W^{\ell-1,p}$ 类算子：步骤 4 建立弗雷德霍姆性质，步骤 5、6 处理 $-1/2\le k-2/q$ 的情形，步骤 7、8 处理一般情形。

\paragraph{以后各步的共同假设。} 从现在起，假定 $D$ 是 $E$ 上的 $W^{\ell-1,p}$ 类\BookRecordIndex{T0210}柯西–黎曼算子，$\ell$ 为正整数、$p>1$、$\ell p>2$；还假定 $k\in\{1,\ldots,\ell\}$、$q>1$ 且 $k-2/q<\ell-2/p$\sourcefootnote{原书此处的共同假设使用严格不等号，但定理 C.2.3 的陈述允许等号；本节后续步骤未单独讨论该端点。这里照录假设，端点情形的证明仍需核实。}，把 $D_F,D_F^*$ 看作从 $W_F^{k,q}$ 到 $W^{k-1,q}$ 的算子。

\paragraph{步骤 4。} 断言 (i) 在一般情形成立。步骤 3、引理~\ref{lem:C-2-2} 和定理~\ref{thm:A-1-5} 给出 $W^{\ell-1,p}$ 类\BookRecordIndex{T0210}柯西–黎曼算子的断言 (i)：引理~\ref{lem:C-2-2} 表明，每个这样的算子都是某个光滑\BookRecordIndex{T0210}柯西–黎曼算子的紧扰动，因此是\BookRecordIndex{T0128}弗雷德霍姆算子，且指标相同。

\paragraph{步骤 5。} 假设 $-1/2\le k-2/q$。则 $D_F,D_F^*$ 满足式~\eqref{eq:C-2-6}。此外，包含 $\ker D_F\hookrightarrow\operatorname{coker}D_F^*$ 及 $\ker D_F^*\hookrightarrow\operatorname{coker}D_F$ 是同构\sourcefootnote{原书步骤 5 此处第二个包含也误写成以 $\operatorname{coker}D_F^*$ 为目标；由伴随配对及本步骤下文可知，$\ker D_F^*$ 应映入 $\operatorname{coker}D_F$。}。

由步骤 4，二者为\BookRecordIndex{T0128}弗雷德霍姆算子，且 $\operatorname{index}(D_F)+\operatorname{index}(D_F^*)=0$。故
\begin{equation}\label{eq:C-2-7}
  \dim\ker D_F+\dim\ker D_F^*
  =\dim\operatorname{coker}D_F+\dim\operatorname{coker}D_F^*.
\end{equation}
由步骤 1(c)，$-1/2\le k-2/q$ 保证 $W^{k,q}$ 自然包含于 $W^{k-1,q}$ 的对偶空间。这样理解后，引理~\ref{lem:C-1-9} 表明：经包含 $W_F^{k,q}(\Sigma,E)\hookrightarrow(W_F^{k-1,q}(\Sigma,E))^*$，$D_F$ 的核映到 $D_F^*$ 的余核（即其像的湮灭子）；同样，$D_F^*$ 的核映到 $D_F$ 的余核。由维数式~\eqref{eq:C-2-7}，这两个包含必为同构。因而对 $W^{\ell-1,p}$ 类\BookRecordIndex{T0210}柯西–黎曼算子，只要 $k\le\ell$ 且 $-1/2\le k-2/q\le\ell-2/p$，式~\eqref{eq:C-2-6} 仍成立。

\paragraph{步骤 6。} 假设 $-1/2\le k-2/q$。则 $D_F,D_F^*$ 满足 (ii) 和 (iii)。

设 $\eta\in L^r(\Sigma,\Lambda^{0,1}T^*\Sigma\otimes_{\mathbb C}E)$、$\xi\in W^{k-1,q}(\Sigma,E)$\sourcefootnote{原书步骤 6 此处将 $\xi$ 的目标丛误印为 $F$；$F$ 只定义在边界上，而式 (C.2.2) 与定理 (ii) 均要求 $\xi$ 是 $E$ 值截面。} 满足式~\eqref{eq:C-2-2}。若 $kq<2$，则 $k=1$，不失一般性可设 $1/r+1/q=1$。（于是 $r$ 是满足 $2/r-1\le k-2/q$ 且 $k=1$ 的最小的、严格大于 $1$ 的实数。）由步骤 5 在目前假设下证明的式~\eqref{eq:C-2-6} 第二式，存在 $\xi_0\in W_F^{k,q}(\Sigma,E)$、$\eta_1\in W_F^{k,q}(\Sigma,\Lambda^{0,1}T^*\Sigma\otimes_{\mathbb C}E)$，使
\begin{equation}\label{eq:C-2-8}
  \xi=\xi_0+D_F^*\eta_1,\qquad D_F\xi_0=0.
\end{equation}
取 $s$ 满足 $1/r+1/s=1$。由步骤 1，有包含 $W^{k-1,q}\hookrightarrow L^s$，故有界线性泛函 $\langle\eta,\cdot\rangle:L^s\to\mathbb R$ 满足
\begin{equation}\label{eq:C-2-9}
  \langle\eta,D_F\zeta\rangle
  =\langle\xi,\zeta\rangle
  =\langle\xi_0+D_F^*\eta_1,\zeta\rangle
  =\langle\xi_0,\zeta\rangle+\langle\eta_1,D_F\zeta\rangle
\end{equation}
（$\zeta\in W_F^{k,q}(\Sigma,E)$）。最后一个等式来自引理~\ref{lem:C-1-9}，以及 $W^{k,q}$ 包含在 $W^{k-1,q}$ 的对偶空间这一事实。令 $\zeta=\xi_0$，得
\[
  \|\xi_0\|_{L^2}^2
  =\langle\xi_0,\xi_0\rangle
  =\langle\eta-\eta_1,D_F\xi_0\rangle=0.
\]
所以 $\xi_0=0$，从而对每个 $\zeta\in W_F^{k,q}(\Sigma,E)$ 有 $\langle\eta-\eta_1,D_F\zeta\rangle=0$。由步骤 5，$\ker D_F^*\hookrightarrow\operatorname{coker}D_F$ 是同构，故 $\eta-\eta_1\in\ker D_F^*$。于是存在 $\eta_0\in W_F^{k,q}(\Sigma,\Lambda^{0,1}T^*\Sigma\otimes_{\mathbb C}E)$，使
\[
  \eta=\eta_0+\eta_1,\qquad D_F^*\eta_0=0.
\]
因此 $\eta\in W_F^{k,q}(\Sigma,\Lambda^{0,1}T^*\Sigma\otimes_{\mathbb C}E)$ 且 $D_F^*\eta=D_F^*\eta_1=\xi$\sourcefootnote{原书步骤 6 末句两处将作用于 $\eta$ 的算子误印为 $D_F$；其定义域是 $E$ 值截面，故此处据类型及式 (C.2.8) 改为 $D_F^*$。}，证明 (ii)。(iii) 由交换 $D_F,D_F^*$ 的角色得出。

\paragraph{步骤 7。} 对每个 $\xi_0\in W_F^{k,q}(\Sigma,E)$，有
\begin{equation}\label{eq:C-2-10}
  D_F\xi_0=0\quad\Longrightarrow\quad\xi_0\in W_F^{\ell,p}(\Sigma,E).
\end{equation}
$D_F^*$ 也一样。此外，包含 $\ker D_F\hookrightarrow\operatorname{coker}D_F^*$ 和 $\ker D_F^*\hookrightarrow\operatorname{coker}D_F$ 是同构，且式~\eqref{eq:C-2-6} 成立。

先设 $kq<2$，取 $\xi_0\in W_F^{k,q}(\Sigma,E)$ 且 $D_F\xi_0=0$。由定理~\ref{thm:B-1-12}，
\[
  \xi_0\in L^r(\Sigma,E),\qquad r:=\frac{2q}{2-kq}.
\]
步骤 1(b) 的证明表明 $2/r-1\le\ell-2/p$ 且有包含 $L^r\hookrightarrow(W^{\ell-1,p})^*$。因此，由引理~\ref{lem:C-1-9}，对每个 $\zeta\in W_F^{\ell,p}(\Sigma,\Lambda^{0,1}T^*\Sigma\otimes_{\mathbb C}E)$，$\langle\xi_0,D_F^*\zeta\rangle=0$。步骤 6 已证明：当 $D_F,D_F^*$ 被视为从 $W_F^{\ell,p}$ 到 $W^{\ell-1,p}$ 的算子时，它们满足定理~\ref{thm:C-2-3}(iii)。在 $\eta=0$ 时应用该断言，得 $\xi_0\in W_F^{\ell,p}(\Sigma,E)$。于是 $kq<2$ 时式~\eqref{eq:C-2-10} 成立；$kq\ge2$ 时论证类似。

由于 $\ell p>2$，有包含 $W^{\ell,p}(\Sigma,E)\hookrightarrow(W^{k-1,q}(\Sigma,E))^*$。由式~\eqref{eq:C-2-10} 和引理~\ref{lem:C-1-9}，$D_F$ 的核自然地包含于 $D_F^*$ 的余核；$D_F^*$ 的核亦然。维数式~\eqref{eq:C-2-7} 表明这些包含是同构。这又表明：当 $-1/2>k-2/q$ 时，式~\eqref{eq:C-2-6} 仍成立。

\paragraph{步骤 8。} (ii) 和 (iii) 在一般情形成立。由步骤 7，$D_F$ 的核包含于 $W^{\ell,p}$，因而包含于 $W^{k-1,q}$ 的对偶空间。重复步骤 6 的论证时只需稍作如下调整。

设 $\eta\in L^r(\Sigma,\Lambda^{0,1}T^*\Sigma\otimes_{\mathbb C}E)$、$\xi\in W^{k-1,q}(\Sigma,E)$ 满足式~\eqref{eq:C-2-2}。由步骤 7 在目前假设下证明的式~\eqref{eq:C-2-6} 第二式，可将 $\xi$ 写成式~\eqref{eq:C-2-8} 的形式；此时 $\xi_0\in W_F^{\ell,p}(\Sigma,E)$，$\eta_1\in W_F^{k,q}(\Sigma,\Lambda^{0,1}T^*\Sigma\otimes_{\mathbb C}E)$，与前面一样。式~\eqref{eq:C-2-9} 对每个 $\zeta\in W_F^{\ell,p}(\Sigma,E)$ 成立。注意现在不再假设 $-1/2\le k-2/q$，所以未必有 $W^{k,q}$ 包含于 $W^{k-1,q}$ 的对偶空间；但 $W^{\ell,p}$ 包含于 $W^{k-1,q}$ 的对偶空间，而且由步骤 1 有 $W^{k,q}$ 包含于 $W^{\ell-1,p}$ 的对偶空间。因此，式~\eqref{eq:C-2-9} 的最后一个等式由引理~\ref{lem:C-1-9} 与这些观察\BookRecordIndex{T0326}推出。与步骤 6 一样，在式~\eqref{eq:C-2-9} 取 $\zeta=\xi_0$，得到 $\xi_0=0$，因而对 $\zeta\in W_F^{\ell,p}(\Sigma,E)$ 有 $\langle\eta-\eta_1,D_F\zeta\rangle=0$。由步骤 7，$\ker D_F^*\hookrightarrow\operatorname{coker}D_F$ 为同构，因此 $\eta-\eta_1\in W_F^{\ell,p}(\Sigma,\Lambda^{0,1}T^*\Sigma\otimes_{\mathbb C}E)$ 且 $D_F^*(\eta-\eta_1)=0$\sourcefootnote{原书步骤 8 此处将作用于 $\eta-\eta_1$ 的算子误印为 $D_F$；该量是 $(0,1)$-形式，故据定义域及前一句改为 $D_F^*$。}。故 $\eta\in W_F^{k,q}(\Sigma,\Lambda^{0,1}T^*\Sigma\otimes_{\mathbb C}E)$ 且 $D_F^*\eta=D_F^*\eta_1=\xi$\sourcefootnote{原书步骤 8 最后一次将作用于 $\eta,\eta_1$ 的算子误印为 $D_F$；据同一类型理由补为 $D_F^*$。}，证明 (ii)。(iii) 仍由交换 $D_F,D_F^*$ 的角色得到。
\end{proof}

\begin{proof}[定理~\ref{thm:C-1-10}(i) 的证明]
设
$\eta\in W^{k-1,q}(\Sigma,\Lambda^{0,1}T^*\Sigma\otimes_{\mathbb C}E)$
满足
\[
  \eta_0\in\ker D_F^*
  \subset W^{\ell,p}(\Sigma,\Lambda^{0,1}T^*\Sigma\otimes_{\mathbb C}E)
  \quad\Longrightarrow\quad
  \langle\eta,\eta_0\rangle=0.
\]
由~\eqref{eq:C-2-6} 和定理~\ref{thm:C-2-3} 证明的步骤 7，
存在 $\xi\in W_F^{k,q}(\Sigma,E)$ 及
$\eta_1\in W^{\ell,p}(\Sigma,\Lambda^{0,1}T^*\Sigma\otimes_{\mathbb C}E)$，
使
\[
  \eta=D_F\xi+\eta_1,\qquad D_F^*\eta_1=0.
\]
在上面的条件中取 $\eta_0=\eta_1$，可得 $\eta_1=0$，
故 $\eta\in\operatorname{im}D_F$。
\end{proof}
